%% defense_mechanism.tex

\section{Defense Mechanism}
\label{sec:defense}

We propose a lightweight, OS-principled defense composed of two mechanisms: \trustorigin{} Tagging (provenance tracking) and the \reauthgate{} (privilege downgrade on recovery). Both are implemented in the \system{} codebase and are evaluated in Section~\ref{sec:evaluation}.

\subsection{Core 1: Data Provenance with \trustorigin{} Tagging}
\label{sec:trustorigin}

\subsubsection{Motivation: From Opaque \code{String} to \code{TaggedContent}}
The root cause of V1 is the system's treatment of error traces as opaque, context-free \code{String} values. When \code{lastErrorTrace} is stored as a bare \code{String}, the provenance of the content---\emph{where it came from}---is lost the moment it is captured. The recovery pipeline therefore has no basis on which to distinguish a trusted kernel-generated error (e.g., a stack overflow in the planner) from an untrusted externally-originated error (e.g., the body of a fetched adversarial web page that triggered a parse failure). The default resolution logic (\code{TrustOrigin.fromMetadata()}) silently fills this gap by assigning \code{SYSTEM\_GENERATED}, which is the most permissive label.

This is a type-safety failure in the provenance dimension: the type system permits untrusted content to flow into a sink that expects trusted content, and no runtime check catches the mismatch. Sanitization does not help, because the trust-bearing signal is not in the content---it is in the \emph{origin} of the content, which the \code{String} representation has discarded.

\subsubsection{Design: Tagged Provenance}
We upgrade the system to carry provenance explicitly by replacing raw \code{String} error traces with a \code{TaggedContent} record that bundles the text with its trust origin. Every content fragment entering the recovery pipeline carries a four-level trust label (Table~\ref{tab:trust-origin}). The label is assigned at the \emph{capture site}---the point where the failure is caught---and propagates with the content through the recovery chain as part of the value's type, not as out-of-band metadata that can be silently dropped.

\begin{lstlisting}[caption={The \code{TaggedContent} record: provenance as part of the value's type. The four \code{TrustOrigin} levels are defined in Table~\ref{tab:trust-origin}.},label={lst:tagged-content}]
public record TaggedContent(
    String text,          // the error / diagnostic content
    TrustOrigin origin,   // trust level (see Table 2)
    String sourceRef      // source identifier, e.g. "web_fetch://evil.com"
) {
    public boolean isTrusted() {
        return origin != TrustOrigin.TOOL_OUTPUT_EXTERNAL;
    }
}
\end{lstlisting}

\begin{table}[htbp]
\centering
\caption{\trustorigin{} trust levels and their handling in recovery.}
\label{tab:trust-origin}
\begin{tabular}{p{0.30\linewidth}p{0.55\linewidth}}
\toprule
\textbf{Trust Level} & \textbf{Recovery Handling} \\
\midrule
\code{SYSTEM\_GENERATED} & High-trust template (\code{[SYSTEM CRITICAL]}). Default for kernel-originated errors. \\
\code{USER\_INPUT} & High-trust template. User is an authorized trust principal. \\
\code{TOOL\_OUTPUT\_INTERNAL} & High-trust template. Internal tools (e.g., \code{bash}) are trusted. \\
\code{TOOL\_OUTPUT\_EXTERNAL} & \textbf{Untrusted template} + ``Do NOT execute'' warning. External tools (e.g., \code{web\_fetch}) process attacker-controlled content. \\
\bottomrule
\end{tabular}
\end{table}

\subsubsection{Implementation}
The \code{TaggedContent} record carries \code{text}, \code{origin}, and \code{sourceRef}. The \code{ToolErrorRecovery} capture site infers the origin from the failing tool: external-content tools (\code{web\_fetch}, \code{web\_scrape}, \code{file\_read} on untrusted paths) tag content as \code{TOOL\_OUTPUT\_EXTERNAL}; internal tools tag as \code{TOOL\_OUTPUT\_INTERNAL}. The \code{ReflectionInjectionRecovery} strategy then selects the prompt template by label (Listing~\ref{lst:trustorigin}).

\begin{lstlisting}[caption={Inferring content origin at the capture site (\code{ToolErrorRecovery.inferOrigin}).},label={lst:capture-site}]
// In ToolErrorRecovery.java
private static final Set<String> EXTERNAL_CONTENT_TOOLS = Set.of(
    "web_fetch", "web_search", "browser_navigate", "browser_snapshot");

private static TrustOrigin inferOrigin(RecoveryContext context) {
    Tool<?> tool = context.originalTool();
    if (tool == null) return TrustOrigin.SYSTEM_GENERATED;
    String name = tool.name();
    return (name != null && EXTERNAL_CONTENT_TOOLS.contains(name))
        ? TrustOrigin.TOOL_OUTPUT_EXTERNAL
        : TrustOrigin.TOOL_OUTPUT_INTERNAL;
}
\end{lstlisting}

\begin{lstlisting}[caption={\trustorigin{}-based template selection in \code{ReflectionInjectionRecovery}.},label={lst:trustorigin}]
TaggedContent tagged = context.taggedError();
String modifier = tagged.isTrusted()
    ? highTrustModifier(lastError, attempt)
    : untrustedModifier(lastError, attempt, tagged.origin());
context.appendPromptModifier(modifier);
\end{lstlisting}

\begin{figure}[htbp]
\centering
\begin{tikzpicture}[
    node distance=0.7cm and 1.0cm,
    box/.style={draw, rounded corners=2pt, minimum height=0.7cm, align=center, font=\footnotesize},
    tool/.style={box, fill=violet!15, draw=violet!60, minimum width=2.8cm},
    capture/.style={box, fill=blue!15, draw=blue!60, minimum width=2.8cm},
    tag/.style={box, fill=cyan!20, draw=cyan!70, minimum width=3.5cm, font=\footnotesize\ttfamily},
    strat/.style={box, fill=orange!20, draw=orange!70, minimum width=3.5cm},
    result/.style={box, fill=green!15, draw=green!60!black, minimum width=3.5cm},
    arr/.style={-Stealth, thick},
    lbl/.style={font=\tiny\itshape, gray!60!black, align=center}
]
\node[tool] (t1) {VFS / Tool\\\code{web\_fetch}};
\node[capture, right=of t1] (t2) {Capture Site\\\code{ToolErrorRecovery}};
\node[tag, right=of t2] (t3) {TaggedContent\\origin=EXTERNAL};
\node[strat, below=1.0cm of t3] (t4) {Recovery Strategy\\\code{ReflectionInjection}};
\node[result, left=of t4] (t5) {Untrusted Template\\+ ``Do NOT execute''};

\draw[arr] (t1) -- (t2) node[lbl, midway, above] {exception};
\draw[arr] (t2) -- (t3) node[lbl, midway, above] {tag};
\draw[arr] (t3) -- (t4) node[lbl, midway, right] {select by \code{origin}};
\draw[arr] (t4) -- (t5) node[lbl, midway, above] {if EXTERNAL};
\end{tikzpicture}
\caption{Data flow with \trustorigin{} tagging. The trust label is assigned once at the capture site and propagates with the content; the recovery strategy consults the label to select the prompt template, isolating externally-originated content in an untrusted framing.}
\label{fig:trustorigin-flow}
\end{figure}

\subsubsection{OS Analogy}
This is the LLM-agent analog of the kernel principle that data copied from userspace retains \code{PAGE\_USER} tagging and may not be executed as kernel code. The trust label is \emph{provenance}, not content---it survives sanitization and propagation.

\subsection{Core 2: Privilege Downgrade \& Unified \reauthgate{}}
\label{sec:reauthgate}

\subsubsection{Motivation: Centralized vs.\ Decentralized Validation}
The root cause of V2 is the \emph{decentralized} validation of recovery actions. In the baseline design, each recovery strategy is responsible for its own permission checks: \code{TopologyMutationStrategy} is expected to call \code{PermissionChecker} before replacing a role, \code{ModelFallbackStrategy} before switching models, and so on. In practice this leads to two failure modes. First, \emph{gaps}: a strategy author who forgets or mis-scopes a check (as in the baseline \code{parseAndValidate(validate=false)} call) leaves a silent hole, because no other component re-checks the action. Second, \emph{inconsistency}: different strategies implement slightly different policies, so the effective permission boundary depends on which strategy the orchestrator happens to dispatch.

This is a classic instance of a \emph{cross-cutting concern} in the software-engineering sense: permission enforcement is an aspect that \emph{every} side-effecting recovery action needs, but it is orthogonal to the strategy's core logic (failure diagnosis and recovery planning). Scattering the concern across strategies couples diagnosis logic to security logic, duplicates enforcement code, and---most importantly---makes the security property \emph{non-local}: there is no single place where one can verify that ``no recovery action bypasses permission checks.''

The principled fix is to centralize the concern at a single choke point. We do this by (i) adding a \code{requiresReauthorization} flag to the \code{RecoveryResult} returned by every strategy, and (ii) intercepting all such results at the \code{RecoveryOrchestrator}---the single component through which every recovery action must flow. The strategies declare \emph{that} an action needs reauthorization; the orchestrator enforces \emph{how}. This separation localizes the security property: the invariant ``no side effect executes without reauthorization'' is now verifiable by inspecting one method rather than eleven strategies.

\begin{figure}[htbp]
\centering
\begin{tikzpicture}[
    node distance=0.7cm and 0.8cm,
    box/.style={draw, rounded corners=2pt, minimum height=0.7cm, minimum width=3.0cm, align=center, font=\footnotesize},
    strat/.style={box, fill=yellow!20, draw=orange!70},
    result/.style={box, fill=orange!20, draw=orange!70, font=\footnotesize\ttfamily},
    gate/.style={box, fill=yellow!30, draw=orange!80, minimum width=3.5cm, font=\footnotesize\bfseries},
    check/.style={box, fill=orange!25, draw=orange!70, minimum width=3.0cm},
    pass/.style={box, fill=green!15, draw=green!60!black},
    block/.style={box, fill=red!15, draw=red!70!black},
    note/.style={box, fill=gray!10, draw=gray!50, font=\tiny\itshape, minimum width=4.5cm},
    arr/.style={-Stealth, thick},
    parr/.style={-Stealth, thick, green!60!black},
    barr/.style={-Stealth, thick, red!70!black},
    lbl/.style={font=\tiny\itshape, gray!60!black, align=center}
]
% Top flow
\node[strat] (s) {Recovery Strategy\\\code{TopologyMutation}};
\node[result, below=of s] (r) {RecoveryResult\\reqReauth=true};
\node[gate, below=of r] (g) {Reauthorization Gate\\\code{RecoveryOrchestrator}};
\node[check, below=of g] (c) {PermissionChecker\\$P_{new} \leq P_{old}$?};

% Branches
\node[pass, below left=0.8cm and 0.3cm of c] (exec) {Execute Side Effect\\resumeNode(newRole)};
\node[block, below right=0.8cm and 0.3cm of c] (deny) {Block\\RecoveryResult.failed};

\node[pass, below=0.6cm of exec] (ok) {Recovery OK\\(privilege preserved)};
\node[block, below=0.6cm of deny] (no) {Attack Stopped};

\draw[arr] (s) -- (r) node[lbl, midway, right] {declare};
\draw[arr] (r) -- (g) node[lbl, midway, right] {intercept};
\draw[arr] (g) -- (c) node[lbl, midway, right] {check};
\draw[parr] (c) -- (exec) node[lbl, midway, above] {pass};
\draw[barr] (c) -- (deny) node[lbl, midway, above] {fail};
\draw[parr] (exec) -- (ok);
\draw[barr] (deny) -- (no);

% Invariant note
\node[note, left=1.2cm of c] (inv) {Privilege Downgrade Invariant:\\$P(\text{action}) \leq P(\text{failed exec})$\\(analog: OS exception handler\\never escalates privileges)};
\draw[dashed, gray!50] (inv) -- (c);
\end{tikzpicture}
\caption{The unified \reauthgate{} architecture. Validation is decoupled from strategy implementation: strategies declare the need for reauthorization via a flag, and a single orchestrator-level gate enforces it before any side effect executes. The privilege-downgrade invariant (left) constrains the checker.}
\label{fig:reauthgate-arch}
\end{figure}

\subsubsection{Principle: Privilege Downgrade Invariant}
The \reauthgate{} enforces the invariant: \emph{no recovery action may exceed the privilege of the failed execution it remedies}. Formally, if a recovery action $a$ remedies a failed execution with privilege $P_{fail}$, then
\[
P(a) \leq P_{fail}.
\]
This mirrors the OS principle that an exception handler executes with the privileges of the faulting context, never higher. For role replacement, the invariant specializes to: the new role's privileges must not exceed the original role's privileges. The \code{PermissionChecker} implements this via \code{checkRoleMutation(originalRole, newRole)}, which verifies that the new role is in the pre-approved downgrade list for the original role (e.g., \code{Python\_Coder} $\rightarrow$ \code{Code\_Reviewer} is allowed; \code{Code\_Reviewer} $\rightarrow$ \code{System\_Admin} is blocked).

\subsubsection{Design: Declare at the Strategy, Enforce at the Orchestrator}
The \code{RecoveryResult} class is extended with a boolean \code{requiresReauthorization} flag (default: \code{false}). Strategies that produce side effects (e.g., role replacement, tool execution) set this flag to \code{true} and store the side effect as a \emph{deferred} callback in the \code{RecoveryContext} metadata, rather than executing it inline (Listing~\ref{lst:defer-side-effect}).

\begin{lstlisting}[caption={Deferring the side effect in \code{TopologyMutationStrategy}; the strategy declares reauthorization need but does not execute.},label={lst:defer-side-effect}]
// In TopologyMutationStrategy.java
public RecoveryResult apply(RecoveryContext context) {
    // ... diagnosis and suggestedRole extraction ...
    context.metadata().put(META_PENDING_SIDE_EFFECT,
        (ctx) -> resumeNode(ctx, suggestedRole)); // deferred, not executed
    return RecoveryResult.okRequiringReauthorization(
        "Topology mutation pending reauth", null); // flag = true
}
\end{lstlisting}

The \code{RecoveryOrchestrator} intercepts every result with \code{requiresReauthorization=true} and routes it through the \reauthgate{} \emph{before} executing the deferred callback (Listing~\ref{lst:reauthgate}). This is a \emph{prevent} control, not a \emph{detect-and-alert} control: a rejected reauthorization blocks the side effect entirely, rather than merely logging it after the fact.

\begin{lstlisting}[caption={\reauthgate{} in \code{RecoveryOrchestrator}: side effect deferred until reauth passes.},label={lst:reauthgate}]
if (result.requiresReauthorization()) {
    ReauthResult reauth = RecoveryReauthorizationGate.check(
        result, context, resolvePermissionChecker());
    if (!reauth.allowed()) {
        return RecoveryResult.failed("denied by reauth gate");
    }
    // reauth passed -> execute deferred side effect
    Function<RecoveryContext, RecoveryResult> se =
        context.metadata().get(META_PENDING_SIDE_EFFECT);
    result = se.apply(context);
}
\end{lstlisting}

\subsubsection{Why This Beats Per-Strategy Checks}
Centralizing reauthorization at the orchestrator yields four concrete advantages over the decentralized baseline, each grounded in established software-engineering practice for managing cross-cutting concerns:
\begin{enumerate}[leftmargin=*]
    \item \textbf{Unified policy.} All security-critical recovery actions traverse the same gate, so the effective policy is identical regardless of which strategy produced the action. This eliminates the inconsistency and gap failure modes of decentralized validation.
    \item \textbf{Prevent, not detect.} Because side effects are deferred until reauthorization passes, an unauthorized action is blocked \emph{before} any state change occurs. A detect-and-alert design (e.g., auditing \code{resumeNode} calls after execution) cannot undo a privilege escalation that has already happened.
    \item \textbf{Separation of concerns.} Recovery strategies focus on failure diagnosis and recovery planning; the security aspect is factored out into a dedicated component. Strategy authors no longer need to be security experts, and security authors no longer need to understand every strategy's internals. This is the same separation that motivates aspect-oriented programming and middleware-based enforcement in classical distributed systems.
    \item \textbf{Open/closed for extension.} A newly added recovery strategy inherits protection automatically by setting \code{requiresReauthorization=true}; it does \emph{not} need to re-implement permission checks. The security property is thus \emph{open for extension} (new strategies are easy to add) but \emph{closed for modification} (the gate code does not change). This scales to the 11+ strategies in \system{}'s orchestrator without per-strategy patches, and means that the security invariant cannot regress when a new strategy is added.
\end{enumerate}

The net effect is that the security property ``no recovery action exceeds the privilege of the failed execution'' becomes a \emph{local, inspectable} property of one orchestrator method, rather than an \emph{emergent} property of eleven independently-authored strategies.

\subsection{Threat Coverage}

\begin{table}[htbp]
\centering
\caption{Defense coverage of attack vectors.}
\label{tab:coverage}
\begin{tabular}{p{0.25\linewidth}p{0.30\linewidth}p{0.30\linewidth}}
\toprule
\textbf{Vector} & \textbf{\trustorigin{}} & \textbf{\reauthgate{}} \\
\midrule
V1: Reflection injection & Template de-escalation (untrusted framing + warning) & N/A (no side effect) \\
V2: Topology mutation & N/A (core dump is internal) & Role whitelist + privilege check (PREVENT) \\
\bottomrule
\end{tabular}
\end{table}

The two mechanisms are complementary: \trustorigin{} addresses content trust (V1), while the \reauthgate{} addresses action privilege (V2). Together they close both attack vectors studied in Section~\ref{sec:introduction}.

\subsection{Robustness Under Unregistered Tools}
\label{sec:unregistered-tools}

A structural limitation of the \trustorigin{} capture site is its reliance on a hardcoded whitelist of external-content tool names (\code{EXTERNAL\_CONTENT\_TOOLS} in Listing~\ref{lst:capture-site}). The current implementation classifies any tool \emph{not} in this set as \code{TOOL\_OUTPUT\_INTERNAL} (high-trust), which is correct for built-in internal tools (e.g., \code{bash}, \code{file\_edit}) but \emph{incorrect} for a newly registered or third-party MCP tool that fetches external content under a non-whitelisted name. Such a tool would silently inherit \code{TOOL\_OUTPUT\_INTERNAL} provenance and be routed to the high-trust recovery template, re-opening V1.

This is a real-world concern: MCP servers can be registered at runtime by user configuration, and their tool names are arbitrary strings. An attacker who can influence the tool-registration configuration---an \emph{extended} threat that is strictly stronger than the baseline external-content-provider model of \S\ref{sec:attacker-model} (where the attacker cannot tamper with system configuration), and a realistic one for multi-tenant deployments---could register a malicious fetch tool under an innocuous name (e.g., \code{docs\_helper}) that is not in \code{EXTERNAL\_CONTENT\_TOOLS}, causing its outputs to be mis-tagged as trusted. We discuss this extension explicitly here rather than in \S\ref{sec:attacker-model} because the baseline threat model already suffices to establish the RCI attack surface; the unregistered-tool concern is a \emph{defense-bypass} scenario for an already-deployed defense, not a new attack vector.

We mitigate this in three layers, ordered by strength:

\begin{enumerate}[leftmargin=*]
    \item \textbf{Fail-closed default (released).} The production code path now treats \emph{unknown} tool names (those not in the explicit internal-tool whitelist) as \code{TOOL\_OUTPUT\_EXTERNAL} by default, flipping the original fail-open behavior. This is the safer default: a mis-tagged internal tool causes only a (recoverable) FAR increase, whereas a mis-tagged external tool causes a (security-critical) ASR increase. The empirical results in Table~\ref{tab:results} were obtained with the original fail-open whitelist, so the defended ASR reported there is a \emph{conservative upper bound} under the released fail-closed default.
    \item \textbf{Explicit trust-class declaration (recommended).} The \code{Tool} interface is extended with a \code{TrustClass trustClass()} method that tool authors must implement at registration time. The capture site reads this declared field instead of inferring from a name list. This eliminates the mismatch-by-omission failure mode entirely, at the cost of requiring every tool to declare its trust class.
    \item \textbf{Static-analysis guard (research prototype).} A compile-time analysis pass inspects each tool's bytecode for external-I/O patterns and flags any tool that performs external I/O but declares \code{TrustClass.INTERNAL}. Concretely, the analysis operates on the ASM-reparsed abstract syntax tree (AST) of each compiled \code{Tool} implementation and applies a three-stage check:
    \begin{enumerate}[leftmargin=*,label=(\alph*)]
        \item \emph{Sink enumeration.} The analysis maintains a registry of \emph{I/O sink} method signatures whose presence in a tool's call graph indicates external-content retrieval: \code{java.net.http.HttpClient.send}, \code{java.net.Socket.connect}, \code{java.lang.ProcessBuilder.start}, \code{java.io.FileInputStream.<init>} (for file-system reads outside a sandboxed root), \code{org.openqa.selenium.WebDriver.get}, and reflective invocations \code{Method.invoke} whose target name matches a known sink pattern. The registry is extensible---new sink classes (e.g., a future MCP client library) can be added without modifying the analysis core.
        \item \emph{Intraprocedural reachability.} For each method in the tool's AST, the analysis performs a forward reachable-in-statement-order scan from the method entry, tracking the set of I/O sinks transitively invoked. The scan is intraprocedural (no inter-procedural call-graph construction), which is conservative: a tool that calls a helper method containing a sink is flagged even if the helper is never reached at runtime. This over-approximation is acceptable because the cost of a false positive is a forced \code{TrustClass.EXTERNAL} declaration (a minor annotation burden), whereas the cost of a false negative is a security-critical mis-tag.
        \item \emph{Taint propagation to return value.} When a sink is reachable, the analysis checks whether the sink's return value (or any value derived from it via a chain of assignments, string concatenations, or collection insertions) flows to the tool's \code{execute()} return. If so, the tool is flagged as \code{EXTERNAL\_CONTENT} regardless of its declared \code{TrustClass}; the build fails with a diagnostic listing the offending sink and the taint path.
    \end{enumerate}
    The analysis catches intentional mis-declaration (e.g., a tool named \code{docs\_helper} that secretly fetches URLs) but is best-effort: a tool that performs external I/O via \code{Method.invoke} with an obfuscated target, or via a side channel (e.g., writing a file that a second tool later reads), can evade it. Strengthening the analysis to inter-procedural call-graph construction and side-channel detection is straightforward in principle but increases build-time cost; we leave the full implementation as future work.
\end{enumerate}

The combination of (i) fail-closed default and (ii) explicit declaration closes the unregistered-tool bypass for any tool whose author is honest; (iii) adds defense-in-depth against dishonest authors. We adopt (i) in the released code; (ii) is implemented but not enforced (existing tools without a declared \code{TrustClass} fall back to the fail-closed default); (iii) is left as future work.

\paragraph{Cross-reference with evaluation.}
The empirical robustness of the defense under unregistered tools is discussed in \S\ref{sec:threats} (Threats to Validity, ``Unregistered-tool degradation''). The released code's fail-closed default is a direct response to the gap identified during this review.
